\documentclass[10pt,letterpaper]{amsart}
\usepackage[margin=0.9in]{geometry}
\usepackage[T1]{fontenc}
\usepackage{lmodern,microtype,amsmath,amssymb,mathtools}
\usepackage[shortlabels]{enumitem}
\usepackage[hidelinks]{hyperref}
\numberwithin{equation}{section}
\newtheorem{theorem}{Theorem}[section]
\newtheorem{proposition}[theorem]{Proposition}
\newtheorem{lemma}[theorem]{Lemma}
\newtheorem{corollary}[theorem]{Corollary}
\theoremstyle{remark}
\newtheorem{remark}[theorem]{Remark}
\newcommand{\cInd}{\operatorname{c\!-\!Ind}}
\newcommand{\End}{\operatorname{End}}
\newcommand{\Hom}{\operatorname{Hom}}
\newcommand{\Irr}{\operatorname{Irr}}
\newcommand{\supp}{\operatorname{supp}}
\newcommand{\C}{\mathbf C}
\setlist[enumerate,1]{label=(\roman*),leftmargin=*}
\setcounter{tocdepth}{1}
\title{Finite algebra data for supercuspidal representations}
\date{1 November 2026}
\hypersetup{pdftitle={Finite algebra data for supercuspidal representations},
 pdfsubject={Jacquet nilpotence, finite matrix parameters, and certified termination}}
\begin{document}
\raggedbottom
\begin{abstract}
We give finite algebra data for irreducible smooth complex supercuspidal
representations of a connected reductive group over a nonarchimedean local field.
At a fixed compact open level and central character, nilpotence relations for
positive Cartan operators define finite-dimensional semisimple Hecke quotients.
A finite boundary calculation determines their multiplication tables. Their simple
modules reconstruct the supercuspidal representations without an additional choice
of an irreducible quotient. A rational certificate for the classical finiteness of
Hecke algebras over their centers gives a terminating search for a uniform stopping
bound. We state the parameters as finite matrices, give their equivalence relation,
and describe the associated group representation by generators and relations.
The construction has no restriction on residue characteristic or splitting of the
group. Algorithmic statements require specified effective local group operations and
exact coefficient arithmetic. The construction generally gives a quotient of a
compact induction; in the sector with nilpotence exponent one it gives a single
irreducible compact induction from a vertex stabilizer.
\end{abstract}
\maketitle
\tableofcontents

\section{The result and the data}
\label{sec:introduction}

Let $F$ be a nonarchimedean local field, let $\mathbf G$ be a connected reductive
$F$-group, and put $G=\mathbf G(F)$ and $Z=Z(G)$. All representations in this
article are smooth complex representations. We fix a decreasing sequence
$(D_e)_{e\ge0}$ of compact open subgroups at one special vertex, as specified
in Section~\ref{sec:cartan-operators}. The sequence is a neighborhood basis
of the identity. The central character will always mean an actual character
$\omega:Z\to\C^\times$, acting by scalars, rather than a generalized central
character.

The result gives data for the supercuspidals with $D_e$-fixed vectors and
central character $\omega$. The data are simple modules of one finite algebra
$Q^{\mathrm{sc}}_{e,\omega}$. This algebra is obtained from double-coset
convolution by imposing powers of explicitly chosen positive Cartan operators.
Its construction does not require a list of supercuspidal representations.

\begin{theorem}\label{thm:main}
For each $e$ and each smooth $\omega$ trivial on $Z\cap D_e$, the construction
in Sections~\ref{sec:cartan-operators}--\ref{sec:certified-level-termination}
gives a finite-dimensional semisimple complex algebra
$Q^{\mathrm{sc}}_{e,\omega}$ with the following properties.
\begin{enumerate}
\item Each simple module $L$ determines an irreducible supercuspidal
representation $\Pi_{e,\omega}(L)$ with
$\Pi_{e,\omega}(L)^{D_e}\simeq L$ and central character $\omega$.
\item Every irreducible supercuspidal representation with these fixed vectors
and central character occurs, and two simple modules give isomorphic group
representations exactly when they are isomorphic.
\item Under the effective input conventions of
Section~\ref{sec:effective-input}, there is a terminating procedure to
compute the multiplication table of $Q^{\mathrm{sc}}_{e,\omega}$ and its
simple modules.
\end{enumerate}
There is no restriction on the residue characteristic or on the splitting
field of $\mathbf G$.
\end{theorem}

Here is the corresponding matrix form of the data. Suppose the computed
algebra is nonzero, and choose a basis $q_1=1,q_2,\ldots,q_n$. Write
$q_iq_j=\sum_k c_{ij}^kq_k$. A datum at $(e,\omega)$ is a positive integer
$d$ and matrices $X_1,\ldots,X_n\in M_d(\C)$ satisfying
\begin{equation}\label{eq:intro-matrix-data}
 X_1=I_d,\qquad X_iX_j=\sum_k c_{ij}^kX_k,\qquad
 \operatorname{span}_{\C}\{X_1,\ldots,X_n\}=M_d(\C).
\end{equation}
The last condition is a finite rank test and is equivalent to simplicity
of the module. Two such data are identified by simultaneous conjugation.
If the algebra is zero, there are no data at that level and central
character. Section~\ref{sec:matrix-data} gives the reconstruction explicitly
as a space of symbols with specified relations and group action.

To parametrize all supercuspidals without repetitions across levels, retain
only data for which averaging over $D_{e-1}$ is zero when $e>0$.
Equivalently, $e$ must be the least level with nonzero fixed vectors.
The central character, least level, and simultaneous-conjugacy class of the
matrices are then complete parameters. One can replace the matrix class by
a primitive central idempotent of $Q^{\mathrm{sc}}_{e,\omega}$.

The proof has three parts. Positive Cartan nilpotence produces a finite
quotient and forces Jacquet modules to vanish. Projectivity of supercuspidals
with fixed central character then makes this quotient semisimple and makes
reconstruction irreducible. Finally, a finite rational certificate for
finiteness over the center supplies an effective stopping bound. The
finiteness theorem used in the last step is classical in characteristic-zero
coefficients; we use the formulation in \cite[Theorem~1.1]{DHKM}.

These are finite algebra parameters. The realization is an explicitly
specified quotient of a compact induction from $ZD_e$. In
Section~\ref{sec:exponent-one}, nilpotence exponent one yields an irreducible
compact induction from a vertex stabilizer. The article does not assert
that every parameter can be converted to a compact inducing pair.

\section{Cartan operators and averaging}
\label{sec:cartan-operators}

Choose a maximal $F$-split torus $S$ and a special vertex $x$ in its reduced
apartment. Let $K$ be the full stabilizer of $x$ in $G$. Thus $Z\subseteq K$
and $K/Z$ is compact. Let $M=C_{\mathbf G}(S)(F)$ be the minimal Levi, and
choose simple relative roots $\alpha_1,\ldots,\alpha_\ell$ and the
corresponding positive system. For any standard parabolic $P=LN$,
write $\overline P=L\overline N$ for its opposite with the same Levi.

We use the standard compact open subgroups at $x$. They can be chosen as a
decreasing neighborhood basis $D_e$, normal in $K$, with simultaneous
Iwahori decompositions for all these parabolics. The required facts are
\cite[Section~I.2, Proposition~7 and Corollaries~8--9]{SS}; a formulation
including the full reduced-building stabilizer is
\cite[Theorem~3.11(a)--(d)]{Solleveld}. We record the precise properties
used below.

\begin{lemma}\label{lem:group-input}
For every $D=D_e$ and every standard parabolic $P=LN$, multiplication gives
an Iwahori decomposition
\[
 D=(D\cap N)(D\cap L)(D\cap\overline N).
\]
The subgroup $D\cap M$ is normal in $M$. An element $a\in M$ whose
translation is zero on the roots of $L$ normalizes $D\cap L$. If its
coordinates contract the positive root groups in $N$, then
\[
 a(D\cap N)a^{-1}\subseteq D\cap N,\qquad
 a^{-1}(D\cap\overline N)a\subseteq D\cap\overline N.
\]
When all these contraction coordinates are strictly positive, the groups
$a^{-n}(D\cap N)a^n$ increase to $N$.
\end{lemma}

\begin{proof}
The cited subgroup construction gives normality and the Iwahori
factorizations. In the notation of \cite[Section~I.2]{SS}, the minimal-Levi
intersection is the good-filtration subgroup $H_{e+}$, normal in
$N_G(S)$, hence normal in $M$.
Conjugation by $M$ shifts the relative root filtrations by its translation
coordinates. Zero coordinates leave the Levi filtrations unchanged;
positive coordinates give the two inclusions. On every root of $N$, strict
positivity makes the filtration index tend to $-\infty$ after conjugation
by $a^{-n}$. The ordered root-group decomposition of $N$ then proves the
exhaustion assertion. These statements use valued root groups, including
multipliable roots, and do not require a split group or a tame extension.
\end{proof}

Write $h:M\to\mathbf R^\ell$ for the simple-root contraction coordinates.
Our sign convention is that $h_i(a)>0$ contracts the corresponding positive
root group. The image $\Lambda=h(M)$ is a lattice, and
$\ker h=M\cap K$. Indeed, $M$ acts on the reduced apartment by translations;
its full translation lattice is commensurable with the image of $X_*(S)$
under evaluation at a uniformizer. The central translation subspace is
spanned by a sublattice, so quotienting it gives a lattice in the reduced
apartment. The simple-root forms are rational on this lattice and give
coordinates on the apartment; see \cite[Lemma~3.3(b)]{Solleveld} and its
building construction.
Put $M^+=h^{-1}(\mathbf R_{\ge0}^\ell)$. Iwasawa and Cartan decomposition give
\begin{equation}\label{eq:group-decompositions}
 G=PK\quad\text{for every standard }P,\qquad G=KM^+K.
\end{equation}
One may use \cite[Theorems~3.13--3.14]{Solleveld} and pass from the extended
stabilizer to $K$; the central translation directions disappear in the
reduced apartment.

Fix $D=D_e$ and a smooth $\omega$ trivial on $Z\cap D$. Let
$\mathcal H_\omega$ be the convolution algebra of locally constant functions
on $G$, compactly supported modulo $Z$, satisfying
$f(zg)=\omega(z)^{-1}f(g)$. Convolution is integration over $G/Z$.
Let $\epsilon$ be the normalized averaging idempotent on $ZD$ with this
central character. It acts as the usual $D$-average $e_D$ on representations
with central character $\omega$. Put
\[
 \mathcal A=\epsilon\mathcal H_\omega\epsilon,
 \qquad [g]=\epsilon\delta_g\epsilon\quad(g\in G),
\]
where $\delta_g$ denotes translation as a convolution multiplier. Thus
$[g]$ acts on $V^D$ as $e_D\pi(g)|_{V^D}$.

\begin{lemma}\label{lem:positive-operators}
For $a,b\in M^+$ and $k_1,k_2\in K$,
\begin{equation}\label{eq:nilpotent-hecke-monoid}
 [a][b]=[ab],\qquad [k_1ak_2]=[k_1][a][k_2].
\end{equation}
If $a$ strictly contracts the radical $N$ of a standard parabolic $LN$
and normalizes $D\cap L$, then, for $v\in V^D$,
\begin{equation}\label{eq:averaging-formula}
 e_D\pi(a^n)v
 =e_{D\cap N}\pi(a^n)v
 =\pi(a^n)e_{a^{-n}(D\cap N)a^n}v\quad(n\ge1).
\end{equation}
\end{lemma}

\begin{proof}
Set $N_D=D\cap N$, $L_D=D\cap L$, and $\overline N_D=D\cap\overline N$.
The Iwahori decomposition gives
$e_D=e_{N_D}e_{L_D}e_{\overline N_D}$. Indeed, $N_DL_D$ is a compact
subgroup and $D=(N_DL_D)\overline N_D$. The pushforward of normalized Haar
measure under this product is normalized Haar measure on $D$: the first
factor acts transitively on $D/\overline N_D$. Apply the same argument to
$N_DL_D$ to obtain the displayed factorization.

The vector $\pi(a^n)v$ is fixed by $L_D$ and $\overline N_D$, by
Lemma~\ref{lem:group-input}. This proves the first equality in
\eqref{eq:averaging-formula}; moving the average across $\pi(a^n)$ proves
the second. Apply the same calculation to the minimal parabolic, allowing
zero contraction coordinates. Since $aN_Da^{-1}\subseteq N_D\subseteq D$,
\[
 e_D\pi(a)e_D\pi(b)v
 =e_D\pi(a)e_{N_D}\pi(b)v
 =e_De_{aN_Da^{-1}}\pi(ab)v
 =e_D\pi(ab)v.
\]
This holds in every smooth representation, hence in the faithful regular
Hecke action. It proves the first identity of
\eqref{eq:nilpotent-hecke-monoid}, without requiring $a$ and $b$ to commute.
Normality of $D$ in $K$ proves the second.
\end{proof}

Choose nonzero Hilbert generators $\gamma_1,\ldots,\gamma_t$ for the monoid
$\Lambda\cap\mathbf R_{\ge0}^\ell$, and choose $b_j\in M^+$ with
$h(b_j)=\gamma_j$. This is a finitely generated monoid because the cone is
rational polyhedral. Every $a\in M^+$ has an ordered factorization
\begin{equation}\label{eq:nilpotent-hecke-cartan}
 a=m_0b_1^{n_1}\cdots b_t^{n_t},\qquad m_0\in M\cap K,\quad n_j\ge0.
\end{equation}
To obtain it, express $h(a)$ in the Hilbert generators and take the ordered
product; the remaining factor lies in $\ker h$. For $N\ge1$ put
\begin{equation}\label{eq:quotient-definition}
 T_j=[b_j],\qquad I_N=\sum_j\mathcal A T_j^N\mathcal A,\qquad
 Q_N=\mathcal A/I_N.
\end{equation}
In relative semisimple rank zero there are no $T_j$, and $Q_N=\mathcal A$.
The quotient is allowed to be zero.

\section{A finite calculation of the quotient}
\label{sec:finite-quotient}


The algebra $\mathcal A$ has at most one basis vector supported on each
$ZD$ double coset. Some cosets may support no vector because of the central
character; they are omitted. Each nonzero $[g]$ is a scalar multiple of the
corresponding basis vector. Equations
\eqref{eq:nilpotent-hecke-monoid} and \eqref{eq:nilpotent-hecke-cartan} show
that $\mathcal A$ is generated by the finite collection
\[
 \mathcal S=\{[k]:k\in K/(ZD)\}\cup\{T_j:1\le j\le t\}.
\]
Choose representatives throughout and discard zero elements. Scalar
central factors in products are evaluated using $\omega$.

Let $C_N$ be the finite set of nonzero double-coset basis vectors supported
in
\begin{equation}\label{eq:nilpotent-hecke-box}
 \bigcup_{0\le n_j<N} K b_1^{n_1}\cdots b_t^{n_t}K.
\end{equation}
Write $V_N$ for their span, $U_N$ for the span of all the omitted basis
vectors, and $p_N:\mathcal A\to V_N$ for the basis projection. Thus
$\mathcal A=V_N\oplus U_N$ as vector spaces. Cartan factorization gives
$U_N\subseteq I_N$. A power $T_j^N$ can nevertheless have retained support
if $N\gamma_j$ has another expression using the Hilbert basis. Accordingly,
\begin{equation}\label{eq:nilpotent-hecke-omitted-ideal}
 I_N=\text{the two-sided ideal generated by }U_N
                       \text{ and }T_1^N,\ldots,T_t^N.
\end{equation}
In particular $Q_N$ is finite-dimensional and
\begin{equation}\label{eq:nilpotent-hecke-dimension-bound}
 \dim Q_N\le [K:ZD]^2N^t.
\end{equation}

The following finite procedure computes the relations among the retained
basis vectors. It does not assume that truncation of convolution is
associative. Let $\mathcal C_N$ be the support set in
\eqref{eq:nilpotent-hecke-box}. For each $s\in\mathcal S$, only finitely
many omitted double cosets meet
\[
 \operatorname{supp}(s)^{-1}\mathcal C_N
 \ \cup\ \mathcal C_N\operatorname{supp}(s)^{-1}.
\]
Both sets are compact modulo $Z$. Compute $p_N(su)$ and $p_N(us)$ for
the omitted basis vectors in these finite sets. Let $R_0\subseteq V_N$
be the span of these boundary vectors together with
$p_N(T_1^N),\ldots,p_N(T_t^N)$. Starting from $R_0$, repeatedly adjoin
\[
 p_N(sR_j)\quad\text{and}\quad p_N(R_js),\qquad s\in\mathcal S.
\]
After at most $\dim V_N$ strict increases, the sequence stabilizes at a
subspace $R_N$.

\begin{lemma}\label{lem:nilpotent-hecke-boundary-closure}
The preceding construction satisfies
\[
 I_N=U_N\oplus R_N,\qquad Q_N=V_N/R_N.
\]
Multiplication in $Q_N$ is induced by $v,w\mapsto p_N(vw)$.
Thus this finite computation retains every relation contributed by double
cosets outside the Cartan box.
\end{lemma}

\begin{proof}
For an omitted basis vector $u$, the coefficient of a retained coset in
$su$ can be nonzero only if its support meets
$\operatorname{supp}(s)^{-1}\mathcal C_N$. The right-product assertion is
the analogous support inclusion. Hence every product of an omitted basis
vector with a generator belongs to $U_N+R_0$. By construction, multiplication
of $R_N$ by any generator, on either side, belongs to $U_N+R_N$.
Since $\mathcal S$ generates $\mathcal A$, the space $U_N+R_N$ is a
two-sided ideal containing $U_N$ and every $T_j^N$. Equation
\eqref{eq:nilpotent-hecke-omitted-ideal} implies $I_N\subseteq U_N+R_N$.

Conversely $U_N\subseteq I_N$. If $u\in U_N$ then $su,us\in I_N$,
and subtracting their omitted-basis components shows that
$p_N(su),p_N(us)\in I_N$. Also $p_N(T_j^N)\in I_N$ because its omitted
component belongs to $U_N\subseteq I_N$. This proves $R_0\subseteq I_N$.
The same argument at every closure step gives $R_N\subseteq I_N$.
The sum is direct because $R_N\subseteq V_N$. The asserted multiplication
now follows from the quotient of the associative algebra $\mathcal A$.
\end{proof}

This is an effective calculation under the usual input assumptions of
effective local group arithmetic: finite double-coset decompositions of
compact sets, multiplication in their finite quotients, and evaluation of
the central character. All boundary products are products of explicitly
given compact sets. Their double-coset lists are finite, and the final
closure is linear algebra with the displayed dimension bound. An abstract
existence theorem for a uniform admissibility bound is not needed.


\section{Reconstruction and supercuspidality}
\label{sec:reconstruction}


For a left $Q_N$-module $L$, form the smooth representation
\[
 \mathcal V(L)=\mathcal H_\omega\epsilon\otimes_{\mathcal A}L.
\]
It is generated by its $D$-fixed vectors, and the unit identity gives
$\mathcal V(L)^D=L$. In general such induction can have a nonzero
subrepresentation with no $D$-fixed vectors. The next argument rules this
out in the present cuspidal situation.

For a parabolic $P=LN$, the unnormalized Jacquet module is the quotient
of $V$ by the span of $\pi(n)v-v$, for $n\in N$ and $v\in V$.
For irreducible smooth complex representations, vanishing of these
quotients for every proper $P$ is equivalent to supercuspidality.

\begin{lemma}\label{lem:nilpotent-hecke-cuspidality}
For every $Q_N$-module $L$, all proper Jacquet modules of $\mathcal V(L)$
vanish.
\end{lemma}

\begin{proof}
For each $j$, let $P_j=M_jN_j$ be the standard parabolic whose Levi roots
vanish on $\gamma_j$, and put $N_{j,0}=D\cap N_j$. The element $b_j$
normalizes $D\cap M_j$: it shifts none of the root filtrations in this
Levi, and $D\cap M$ is normal in $M$. It strictly contracts $N_j$ and strictly
expands its opposite radical. Thus Lemma~\ref{lem:positive-operators} gives, for $v\in\mathcal V(L)^D$,
\[
 T_j^Nv=e_D\pi(b_j^N)v
       =\pi(b_j^N)e_{b_j^{-N}N_{j,0}b_j^N}v.
\]
The left side is zero by the definition of $Q_N$. Hence the image of $v$
in the $N_j$-coinvariants is zero. Iwasawa decomposition $G=P_jK$ and
normality of $D$ in $K$ imply that these images generate the entire Jacquet
module: $\mathcal V(L)$ is generated by its $D$-fixed vectors, and $K$
preserves that space. Thus the Jacquet module is zero. The Hilbert basis
contains the primitive lattice generator on every extreme ray: a point
on such a ray can only be decomposed as a sum of points on the same ray.
The parabolics $P_j$ therefore include every maximal standard parabolic.
Conjugation and transitivity of Jacquet functors give the assertion for
every proper parabolic. In relative rank zero there are no proper
parabolics and the assertion is immediate.
\end{proof}

\begin{lemma}\label{lem:projectivity}
An irreducible supercuspidal representation is projective in the category
of smooth representations with its actual central character fixed.
\end{lemma}

\begin{proof}
We recall the coefficient proof of \cite[Theorem~5.4.1]{Casselman}.
First suppose the central character is unitary. A supercuspidal is
unitarizable: integrate the squared absolute value of a nonzero matrix
coefficient over $G/Z$ to obtain a positive invariant form. The integral
is finite because supercuspidal matrix coefficients have compact support
modulo $Z$. For a unit vector $v$, with the inner product linear in its first argument,
Schur orthogonality gives the idempotent
$f(g)=d_\pi\langle\pi(g^{-1})v,v\rangle$ in $\mathcal H_\omega$, where
$d_\pi$ is the formal degree. Its left ideal $\mathcal H_\omega f$ is
isomorphic to the representation. Consequently,
\[
 \Hom_{\mathcal H_\omega}(\mathcal H_\omega f,V)\simeq fV.
\]
Multiplication by an idempotent is an exact functor, so this module is
projective. For general $\omega$, twist by a positive unramified character
whose restriction to $Z$ is $|\omega|^{-1}$. Such an extension exists on
the real character spaces of the central valuation lattices; see
\cite[Lemma~5.2.5]{Casselman}. Twisting
is an exact equivalence and transfers projectivity back. This is the
usual fixed-central-character projectivity argument; it uses no
compact-induction assertion.
\end{proof}

\begin{theorem}\label{thm:nilpotent-hecke-classification}
The algebra $Q_N$ is finite-dimensional and semisimple. Its simple modules
are in bijection with the irreducible supercuspidal representations $\pi$
having central character $\omega$, nonzero $D$-fixed vectors, and
\[
 (e_D\pi(b_j)|_{\pi^D})^N=0\qquad(1\le j\le t).
\]
The two maps are $L\mapsto\mathcal V(L)$ and $\pi\mapsto\pi^D$.
Every irreducible supercuspidal with central character $\omega$ and
$\pi^D\ne0$ occurs for some $N$, and one may take $N=\dim\pi^D$.
\end{theorem}

\begin{proof}
First take $L=Q_N$ with its regular left action. The representation
$\mathcal V(Q_N)$ is generated by the finite-dimensional space $Q_N$
and is cuspidal by the lemma. Every nonzero finitely generated smooth
representation has an irreducible quotient: a maximal proper submodule
exists because the union of a chain of proper submodules cannot contain
a fixed finite generating set. Any irreducible quotient of
$\mathcal V(Q_N)$ is supercuspidal, by exactness of Jacquet functors.
Projectivity splits off that quotient.

The complement is still generated by its $D$-fixed vectors, because it is
the image of a $G$-equivariant projection of a $D$-generated representation.
Every nonzero quotient of a $D$-generated representation has nonzero
$D$-fixed vectors. Thus each splitting strictly decreases the dimension
of the remaining $D$-fixed space. After at most $\dim Q_N$ splittings,
the complement is zero. This proves that $\mathcal V(Q_N)$ is a finite
direct sum of irreducible supercuspidals.

For an irreducible smooth $\pi$ with $\pi^D\ne0$, its $D$-fixed space is
a simple $\mathcal A$-module. Indeed, a nonzero $\mathcal A$-submodule
generates $\pi$ under $G$, and averaging that generated space by $e_D$
recovers exactly the original submodule. Taking $D$-fixed vectors in the
preceding direct sum therefore makes the regular module $Q_N$ a direct
sum of simple $Q_N$-modules. Hence $Q_N$ is semisimple.

For a simple $Q_N$-module $L$, the same splitting argument applied to
$\mathcal V(L)$ gives a finite supercuspidal direct sum. Its $D$-fixed
space is the simple module $L$, so the sum has exactly one irreducible
summand, with multiplicity one. Conversely, if $\pi$ satisfies the
displayed nilpotence conditions, its fixed space is a simple $Q_N$-module.
The natural nonzero map $\mathcal V(\pi^D)\to\pi$ is then an isomorphism.

Finally, for a supercuspidal $\pi$ the Jacquet averaging argument makes
each $T_j$ locally nilpotent on $\pi^D$: the increasing compact groups
$b_j^{-n}N_{j,0}b_j^n$ exhaust $N_j$, and a vector with zero image in
the Jacquet module is killed by averaging over a sufficiently large one
of these groups. This space is finite-dimensional
by admissibility. A nilpotent operator on a space of dimension $d$ has
$d$th power zero, proving the last assertion.
\end{proof}

The index $N$ bounds nilpotence exponents, not the dimension of a simple
module. A simple $Q_N$-module may have dimension larger than $N$.
The inclusions $I_{N+1}\subseteq I_N$ give surjections $Q_{N+1}\to Q_N$,
and the union of their simple spectra gives exactly the supercuspidal
representations with these fixed $D$ and $\omega$. The construction does
not by itself provide an effective stabilization index for this sequence.
Section~\ref{sec:certified-level-termination} supplies such an index by
a separate finite certificate for the finiteness theorem for Hecke algebras.


\section{A certified stopping bound}
\label{sec:certified-level-termination}

Write $\epsilon_e$, $\mathcal A_{e,\omega}$, and $Q_{e,\omega,N}$ for
the preceding idempotent, corner algebra, and quotient formed using
$D=D_e$ and the central character $\omega$. We have computed the quotient
for each specified $N$. We now remove the search for a sufficiently large $N$. The bound
will depend on $G$, the chosen subgroup $D_e$, and the fixed conventions
for local group computations, but not on $\omega$ or on a supercuspidal
representation. Its existence uses the classical finiteness of a Hecke
algebra over its center. We give a finite certificate for this finiteness,
so that an effective search for the certificate has a proved stopping
condition. This is an application of the finiteness theorem, not a new
proof of that theorem.

For this purpose, work first without fixing a central character. Put
$\mathcal H_D=\mathbf Q[D\backslash G/D]$, with Haar measure normalized
by $\operatorname{vol}(D)=1$. Its double-coset basis has rational
structure constants, and each product has finite support. The effective
input consists of an enumeration of these double cosets, their finite
convolution products, and the Cartan and finite quotient operations used
above. These are requirements on an encoding of the input, rather than
restrictions on residue characteristic, ramification, or the group.

There is a finite, explicitly specified set $\mathcal S_D$ of algebra
generators. Indeed, choose representatives of $K/(ZD)$, generators and
their inverses of the finitely generated abelian group $Z/(Z\cap D)$,
and the positive Cartan elements $b_j$. Take the corresponding compressed
operators. An element of $K$ is a product of a central element, one of the
chosen representatives, and an element of $D$. The Cartan decomposition
and the positive multiplication identity
\eqref{eq:nilpotent-hecke-monoid}, applied before specialization of the
central character, then express every double-coset basis vector in the
algebra generated by $\mathcal S_D$. The finite generation of
$Z/(Z\cap D)$ follows from the compactness of the maximal bounded
subgroup of the center and the finite rank of its valuation lattice.

We use the following algebraic observation. All elements appearing in a
certificate have finite support in the specified algebra basis.

\begin{lemma}\label{lem:finite-central-certificate}
Let $A_0$ be a unital rational algebra with an effective countable basis,
effective finite products, and a specified finite generating set $S$.
Suppose $A_0\otimes_{\mathbf Q}E$ is finite over its center for some
field extension $E/\mathbf Q$. There is a terminating search for elements
$z_1,\ldots,z_r,a_1,\ldots,a_m$ in $A_0$ and polynomials
$f_{s,j,k}\in\mathbf Q[X_1,\ldots,X_r]$ such that
\begin{equation}\label{eq:finite-central-certificate}
 a_1=1,\qquad [z_i,s]=0,\qquad
 sa_j=\sum_{k=1}^{m}f_{s,j,k}(z_1,\ldots,z_r)a_k
 \quad(s\in S).
\end{equation}
For every such certificate, every finite-dimensional irreducible complex
$A_0$-module $V$ satisfies $(\dim V)^2\le m$.
\end{lemma}

\begin{proof}
Since $S$ generates $A_0$, its center is the kernel of the rational linear map
\[
 A_0\longrightarrow\bigoplus_{s\in S}A_0,
 \qquad a\longmapsto([a,s])_{s\in S}.
\]
The set $S$ is finite, so flat base change gives
$Z(A_0)\otimes E=Z(A_0\otimes E)$. Choose finitely many module generators
of $A_0\otimes E$ over this center, and write each as a finite sum of
rational algebra elements with coefficients in $E$. Collect these
rational elements as $a_1,\ldots,a_m$, adding $a_1=1$ if necessary.
The module $M_0=\sum_j Z(A_0)a_j$ satisfies $M_0\otimes E=A_0\otimes E$.
Faithfulness of extension of scalars implies $M_0=A_0$.

For each $s\in S$ and each $j$, express $sa_j$ in these module generators
with coefficients in $Z(A_0)$. There are finitely many such coefficients;
list them as $z_1,\ldots,z_r$. This supplies a rational certificate, with
polynomials of degree at most one. It is unnecessary to compute the
entire center first.

Each displayed condition is a finite equality in the specified basis.
Enumerate all finite lists of finite-support elements with rational
coefficients and all finite lists of rational polynomials. Check the
identities by exact convolution and rational arithmetic. An enumeration
by finite descriptions ensures that every candidate is eventually
examined. The existence just proved ensures that the search terminates.

Conversely, a passing certificate suffices on its own. The equations
$[z_i,s]=0$ make every $z_i$ central, since $S$ generates the algebra.
The module
\[
 M=\sum_{k=1}^{m}\mathbf C[z_1,\ldots,z_r]a_k
\]
contains $1$ and is stable under left multiplication by every $s\in S$.
Induction on word length gives $M=A_0\otimes\mathbf C$.
On an irreducible finite-dimensional module $V$, each $z_i$ acts by a
scalar. The image of the algebra in $\operatorname{End}_{\mathbf C}(V)$
is therefore spanned by the images of $a_1,\ldots,a_m$. Burnside's theorem
identifies this image with the full endomorphism algebra. Its dimension
is $(\dim V)^2$, which proves the bound.
\end{proof}

For $A_0=\mathcal H_D$, apply
\cite[Theorem~1.1]{DHKM} with $R=\mathbf Q_\ell$, where
$\ell\ne p$. Its hypotheses hold for every compact open subgroup over
every nonarchimedean local field. Thus $\mathcal H_D\otimes\mathbf Q_\ell$
is finite over its center, as required by the lemma. The certificate
itself is rational; no computation with arbitrary $\ell$-adic numbers
is needed.

Fix one deterministic enumeration of the certificates, and let $m_e$ be
the number of module generators in its first passing certificate for
$D_e$. Define
\[
 B_e=\max\{1,\lfloor\sqrt{m_e}\rfloor\}.
\]
This integer is computed before a central character is supplied.
Irreducible smooth complex representations are admissible; see, for example,
\cite[Proposition~4.10(iv)]{HV} for arbitrary local field characteristic.
Admissibility and the usual fixed-vector simplicity imply that
$\pi^{D_e}$ is a finite-dimensional simple module of the ordinary
Hecke algebra whenever $\pi$ is irreducible and $\pi^{D_e}\ne0$.
Lemma~\ref{lem:finite-central-certificate} therefore gives
\begin{equation}\label{eq:certified-fixed-dimension-bound}
 \dim\pi^{D_e}\le B_e.
\end{equation}

\begin{theorem}\label{thm:certified-cuspidal-stabilization}
For every smooth central character $\omega$ trivial on $Z\cap D_e$,
put $Q^{\mathrm{sc}}_{e,\omega}=Q_{e,\omega,B_e}$. Its simple modules
give all irreducible supercuspidal representations with central character
$\omega$ and nonzero $D_e$-fixed vectors. For every $N\ge B_e$, the
canonical surjection
\[
 Q_{e,\omega,N}\longrightarrow Q^{\mathrm{sc}}_{e,\omega}
\]
is an isomorphism. Computing $B_e$, forming this finite algebra, and
decomposing it is a terminating procedure under the effective input
conventions already stated.
\end{theorem}

\begin{proof}
If $\pi$ is supercuspidal with nonzero $D_e$-fixed vectors, every positive
Cartan operator on $\pi^{D_e}$ is nilpotent. Its nilpotence index is at
most $\dim\pi^{D_e}\le B_e$. Thus Theorem~\ref{thm:nilpotent-hecke-classification} puts every such $\pi$
in the simple spectrum of $Q^{\mathrm{sc}}_{e,\omega}$.

For $N\ge B_e$, take any simple $Q_{e,\omega,N}$-module $L$.
The same theorem realizes it as $\pi^{D_e}$ for a supercuspidal $\pi$.
The dimension bound implies $T_j^{B_e}L=0$ for every $j$, so $L$
factors through $Q^{\mathrm{sc}}_{e,\omega}$. The kernel of the
displayed surjection therefore annihilates every simple module of
$Q_{e,\omega,N}$. That algebra is semisimple, hence the kernel is zero.
The certificate search terminates by
Lemma~\ref{lem:finite-central-certificate}; the boundary calculation
terminates by Lemma~\ref{lem:nilpotent-hecke-boundary-closure}.
Decomposition of the resulting finite semisimple algebra completes the
procedure. The coefficient interface for the specified central character
includes exact arithmetic, decidable equality, and algebraic extensions
needed to split finitely many polynomials. This is separate from the
certificate search, which uses rational arithmetic only.
\end{proof}

An equality between two successive quotients, by itself, is not a valid
stopping certificate. For example, let $A=\mathbf C\oplus M_r(\mathbf C)$,
let $T=(0,J_r)$ with $J_r$ a nilpotent Jordan block, and take $r\ge3$.
Then $A/(AT^NA)=\mathbf C$ for $1\le N<r$, whereas $A/(AT^rA)=A$.
This finite algebra example concerns the proposed stopping rule, not an
assertion that these matrices occur in a reductive-group Hecke algebra.
The bound $B_e$ avoids reliance on an unverified plateau.


\section{Parameters and the group action}
\label{sec:matrix-data}

\begin{corollary}\label{cor:certified-complete-algebra-data}
Fix the vertex, filtration, Cartan representatives, and certificate
enumeration. All irreducible supercuspidal representations are parametrized
by triples $(e,\omega,c)$, where $e\ge0$, $\omega$ is smooth and trivial
on $Z\cap D_e$, and $c$ is a primitive central idempotent of the
explicitly computed finite algebra $Q^{\mathrm{sc}}_{e,\omega}$.
When $e>0$, require
\[
 u_{e-1}c=0,
\]
where $u_{e-1}$ is the averaging idempotent of $D_{e-1}$, viewed in the
$D_e$-corner and then in $Q^{\mathrm{sc}}_{e,\omega}$.
Two such triples give isomorphic representations if and only if they
are equal.
\end{corollary}

\begin{proof}
Each primitive central idempotent $c$ cuts out a simple matrix algebra
$Q^{\mathrm{sc}}_{e,\omega}c$, with exactly one isomorphism class of
simple left modules. Choose a primitive idempotent $p$ in that summand
and put $L=Q^{\mathrm{sc}}_{e,\omega}p$. The corresponding representation
is
\[
 \Pi(e,\omega,c)
 =\mathcal H_\omega\epsilon_e
      \otimes_{\epsilon_e\mathcal H_\omega\epsilon_e}L.
\]
Different choices of $p$ give isomorphic simple modules, and hence
isomorphic representations. The preceding theorem and
Theorem~\ref{thm:nilpotent-hecke-classification} prove irreducibility,
supercuspidality, and exhaustion at the chosen level.

The representation of the matrix summand on $L$ is faithful. Consequently
$u_{e-1}c=0$ is equivalent to $u_{e-1}L=0$, and hence to the vanishing
of $\Pi(e,\omega,c)^{D_{e-1}}$. Since the $D_e$ form a decreasing
neighborhood basis, this makes $e$ exactly the least level with nonzero
invariants. The central character and this least level are intrinsic.
At equal $e$ and $\omega$, distinct primitive central idempotents give
nonisomorphic simple modules, proving the asserted equivalence relation.
Conversely, for any supercuspidal $\pi$, take its least level $e$, its
central character $\omega$, and the central idempotent supporting the
simple module $\pi^{D_e}$. These form the required triple.
\end{proof}



\subsection{Finite matrix data}

Fix $(e,\omega)$, put $Q=Q^{\mathrm{sc}}_{e,\omega}$, and assume $Q\ne0$.
Choose a basis $q_1=1,q_2,\ldots,q_n$ and its multiplication constants
$c_{ij}^k$. When $e>0$, compute $u_{e-1}=\sum_i\eta_iq_i$ in this basis.
Explicitly, before passing to $Q$, this element is
\[
 u_{e-1}=\frac{1}{[D_{e-1}:D_e]}
             \sum_{h\in D_{e-1}/D_e}[h].
\]
If $\omega$ is nontrivial on $Z\cap D_{e-1}$, this average is zero in
the central-character algebra. No extra restriction on that intersection
is imposed on the data.

\begin{theorem}\label{thm:matrix-parameters}
The irreducible supercuspidal representations are parametrized by data
\[
 \mathcal D=(e,\omega,d;X_1,\ldots,X_n),
\]
where $1\le d\le B_e$, the matrices satisfy
\eqref{eq:intro-matrix-data}, and, when $e>0$, they also satisfy
$\sum_i\eta_iX_i=0$. The associated representation is
\begin{equation}\label{eq:matrix-realization}
 \Pi(\mathcal D)=\mathcal H_\omega\epsilon_e
       \otimes_{\mathcal A_{e,\omega}}\C^d,
 \qquad q_i v=X_i v.
\end{equation}
Two data are equivalent exactly when their $e$, $\omega$, and $d$ agree
and an invertible matrix $S$ satisfies $SX_i=X_i'S$ for all $i$.
For two valid data with the same $e$ and $\omega$, it is enough to find a
nonzero simultaneous intertwiner; this is a finite system of linear equations.
\end{theorem}

\begin{proof}
The multiplication equations and $X_1=I_d$ define a unital representation
of $Q$ on $\C^d$. Its image is precisely the linear span of the $X_i$,
because the $q_i$ form a basis of the algebra. Burnside's theorem says that
this image is $M_d(\C)$ exactly when the module is simple. Thus the rank
condition in \eqref{eq:intro-matrix-data} is exactly the required
irreducibility condition.

Theorem~\ref{thm:certified-cuspidal-stabilization} and
Theorem~\ref{thm:nilpotent-hecke-classification} now prove that
\eqref{eq:matrix-realization} is irreducible supercuspidal, and that every
supercuspidal at this level is obtained. Its fixed-vector space is $\C^d$,
so \eqref{eq:certified-fixed-dimension-bound} gives $d\le B_e$.
The equation $\sum_i\eta_iX_i=0$ says that the preceding level has no fixed
vectors. It therefore enforces the least level, as in
Corollary~\ref{cor:certified-complete-algebra-data}. Isomorphisms of the
simple modules are precisely simultaneous changes of basis. A nonzero
intertwiner between simple modules has zero kernel and full image, hence
is an isomorphism. This proves the last assertion and exhaustion without
repetitions across levels.
\end{proof}

\subsection{The representation by generators and relations}

We spell out the group representation in \eqref{eq:matrix-realization}.
This also specifies an action for every element of $G$, rather than only
for the finite Hecke generators. Put $D=D_e$ and $L=\C^d$. Choose elements
$s_1,\ldots,s_r\in G$ such that
$a_i=\epsilon_e\delta_{s_i}\epsilon_e$ generate $\mathcal A_{e,\omega}$;
one may use the representatives in $K/(ZD)$ and the $b_j$. Let $A_i$ be
the matrix of the image of $a_i$ on $L$, obtained from the matrices $X_k$
and the computed coordinates of that image in $Q$.

Start with symbols $[g,v]$, for $g\in G$ and $v\in L$, linear in $v$,
and form their finite linear combinations. Impose
\begin{equation}\label{eq:symbol-central}
 [gd,v]=[g,v],\qquad [gz,v]=\omega(z)[g,v]
 \quad(d\in D,\ z\in Z).
\end{equation}
This space is $\cInd_{ZD}^G(\omega\otimes L)$, where
$zd$ acts on $L$ as the scalar $\omega(z)$. Choose representatives
$d_{i,1},\ldots,d_{i,m_i}$ of
$D/(D\cap s_iDs_i^{-1})$, with
$m_i=[D:D\cap s_iDs_i^{-1}]$. Impose in addition
\begin{equation}\label{eq:symbol-hecke}
 \frac1{m_i}\sum_{j=1}^{m_i}[g d_{i,j}s_i,v]=[g,A_i v]
 \quad(g\in G,\ v\in L,\ 1\le i\le r).
\end{equation}
The group action on the resulting quotient is
\begin{equation}\label{eq:symbol-action}
 h[g,v]=[hg,v]\quad(h\in G).
\end{equation}
All vectors are finite sums; there is no completion.

\begin{proposition}\label{prop:symbol-realization}
The quotient defined by \eqref{eq:symbol-central} and
\eqref{eq:symbol-hecke}, with action \eqref{eq:symbol-action}, is canonically
isomorphic to \eqref{eq:matrix-realization}. For every valid datum it is
nonzero, irreducible, and supercuspidal. The map $v\mapsto[1,v]$ identifies
$L$ with its $D$-fixed space.
\end{proposition}

\begin{proof}
Write $u_g=\delta_g*\epsilon_e\in\mathcal H_\omega\epsilon_e$. These
elements satisfy $u_{gd}=u_g$ and $u_{gz}=\omega(z)u_g$. They form the
usual basis with central-character relations for compact induction.
Normalized averaging gives
\[
 u_g*a_i=\frac1{m_i}\sum_{j=1}^{m_i}u_{g d_{i,j}s_i}.
\]
Thus \eqref{eq:symbol-hecke} imposes exactly the balancing relation
$(u*a_i)\otimes v=u\otimes A_i v$. Relations for the finite generating
set imply balancing for a word in the generators by induction on its
length, and then for every element of $\mathcal A_{e,\omega}$ by linearity.
The quotient is consequently the tensor product in
\eqref{eq:matrix-realization}.

Averaging this tensor product by $\epsilon_e$ gives
$\mathcal A_{e,\omega}\otimes_{\mathcal A_{e,\omega}}L=L$, with the stated
map. Smoothness also follows directly: a finite sum of symbols is fixed
by an intersection of finitely many conjugates of $D$. The remaining
claims follow from Theorem~\ref{thm:nilpotent-hecke-classification}.
\end{proof}

\section{The computation and its input}
\label{sec:effective-input}

The algebraic assertions above concern every connected reductive $G$ over
a nonarchimedean local field and every smooth complex central character.
For a literal algorithm, we specify the following effective input.
\begin{enumerate}
\item The chosen subgroups, Cartan lattice and Hilbert generators are
available. Ordinary $D$ double cosets can be enumerated and their finite
convolution products can be computed exactly. Finite double-coset lists
for the compact sets occurring in the boundary calculation are available,
with the corresponding central factors when working modulo $Z$.
\item The prescribed central character can be evaluated on the central
factors in these operations. The coefficient field containing these
values has exact arithmetic and decidable equality, with effective
algebraic extensions and factorization sufficient to split the finite
algebras that occur.
\end{enumerate}
The character is part of the input. In particular, arbitrary unspecified
complex numbers are not being treated as computable numbers, and the
procedure does not enumerate all complex central characters.

For a specified $(e,\omega)$, the procedure is as follows.
\begin{enumerate}
\item Search the rational certificates of
Lemma~\ref{lem:finite-central-certificate} for the ordinary Hecke algebra
at $D_e$. The first passing certificate gives $m_e$ and $B_e$.
\item Put $N=B_e$. List the retained double cosets in
\eqref{eq:nilpotent-hecke-box} and compute the boundary relations and the
subspace $R_N$ of Lemma~\ref{lem:nilpotent-hecke-boundary-closure}.
\item Form the multiplication table of $Q=V_N/R_N$. If $Q=0$, return the
empty list. Otherwise decompose the finite semisimple algebra, choose a
simple module for each matrix factor, and compute the matrices of its basis.
\item For parameters of least level, when $e>0$, discard the modules on
which $u_{e-1}$ acts nontrivially. Retain one module per matrix factor and use
\eqref{eq:symbol-central}--\eqref{eq:symbol-action} for the associated group
representation.
\end{enumerate}
The certificate search terminates by finiteness over the center. The
boundary closure makes at most $\dim V_N$ strict increases, and the
remaining operations are finite algebra calculations. This proves the
algorithmic assertion of Theorem~\ref{thm:main}; the reconstruction and
exhaustion assertions follow from
Theorem~\ref{thm:certified-cuspidal-stabilization} and
Theorem~\ref{thm:matrix-parameters}. No complexity bound is asserted for
the exhaustive certificate search.

\section{Compact induction and examples}
\label{sec:exponent-one}

\begin{proposition}\label{prop:nilpotent-hecke-exponent-one}
Every simple $Q_1$-module $W$ is an irreducible representation of $K$,
trivial on $D$ and with central character $\omega$, and its associated
supercuspidal is
\[
 \mathcal V(W)\simeq\operatorname{c-Ind}_K^G W.
\]
Consequently, the supercuspidals with nonzero $D$-fixed vectors on which
every $T_j$ acts by zero are classified by these irreducible compact
inducing data. This assertion holds for every connected reductive group
and every residue characteristic.
\end{proposition}

\begin{proof}
For $N=1$, the support set in \eqref{eq:nilpotent-hecke-box} is $K$.
Thus $Q_1$ is a quotient of the finite twisted group algebra of $K/(ZD)$
with central character $\omega$. Its simple module $W$ is irreducible
as a $K$-representation. By
Theorem~\ref{thm:nilpotent-hecke-classification}, the associated
$\pi=\mathcal V(W)$ is irreducible supercuspidal and $\pi^D=W$.
Moreover,
\[
 e_D\pi(g)|_W=0\qquad(g\notin K),
\]
because every omitted double-coset basis vector is zero in $Q_1$.

Twist by a positive unramified character to make the central character
unitary. This character is trivial on the compact group $D$, so the
vanishing assertion and the identification of $D$-fixed vectors are
unchanged. Equip $\pi$ with an invariant positive definite Hermitian
form. The operator $e_D$ is the orthogonal projection onto $W$.
If $gK\ne hK$, then for $v,w\in W$ one has
\[
 \langle\pi(g)v,\pi(h)w\rangle
 =\langle v,e_D\pi(g^{-1}h)w\rangle=0.
\]
Hence the spaces $\pi(g)W$, indexed by $G/K$, form an algebraic direct
sum. They span $\pi$ because $W$ generates it. The natural map from
$\operatorname{c-Ind}_K^G W$ to $\pi$ is therefore both injective and
surjective. Twisting back proves the assertion for the original central
character.
\end{proof}



\subsection{Groups compact modulo their center}

If $\mathbf G$ has relative semisimple rank zero, its reduced building is
a point, $K=G$, and there are no positive Cartan operators. The algebra
$Q_N=\mathcal A$ is the finite twisted group algebra of $G/(ZD)$ determined
by $\omega$. Its simple modules are already representations of $G$,
trivial on $D$. There are no proper parabolic subgroups, so they are all
supercuspidal. In particular, if $\mathbf G$ is a torus, then $Z=G$ and
$Q_N=\C$ for each admissible $\omega$. The sole simple module reconstructs
the one-dimensional representation $\omega$; the least-level condition
records the first $D_e$ contained in its kernel.

\subsection{A concrete sector in residue characteristic two}

Let $F$ have residue field $\mathbf F_2$, let $G=\operatorname{PGL}_2(F)$,
and take $K=\operatorname{PGL}_2(\mathcal O_F)$ and
$D=\ker(K\to\operatorname{PGL}_2(\mathbf F_2))$. Put
$b=\operatorname{diag}(\varpi,1)$, viewed in $G$. Then $K/D\simeq S_3$,
and there is one positive Cartan generator $T=[b]$.

\begin{proposition}\label{prop:dyadic-example}
For these choices, $Q_1=\mathcal A/(\mathcal A T\mathcal A)\simeq\C$.
Its simple module is the sign character of $S_3$, inflated to $K$, and
its associated representation is $\cInd_K^G\operatorname{sgn}$.
Thus the matrix data in this sector are scalar signs on $K/D$ and the
scalar zero on $T$.
\end{proposition}

\begin{proof}
The algebra $Q_1$ is a semisimple quotient of $\C[S_3]$. If $W$ is one
of its simple modules, Proposition~\ref{prop:nilpotent-hecke-exponent-one}
says that $\cInd_K^G W$ is supercuspidal. Let $N$ be a maximal unipotent
subgroup with integral subgroup $U=K\cap N$. Its orbit through the
identity coset in $G/K$ contributes the summand $\cInd_U^N W$ to the
restriction of this compact induction to $N$. Taking $N$-coinvariants
of that summand gives $W_U$. Since the whole Jacquet module is zero,
$W_U=0$. The image of $U$ in $S_3$ is generated by a transposition.
Among the trivial, sign, and two-dimensional standard representations
of $S_3$, only sign has zero invariants under that subgroup. In
characteristic zero invariants and coinvariants for a finite group have
the same dimension. Thus sign is the only possible simple module of $Q_1$.

It remains to show that it occurs. Write $a_n=\operatorname{diag}(\varpi^n,1)$
for $n\ge1$, and let $U^+$ and $U^-$ be the upper and lower integral
unipotent subgroups. Direct multiplication gives
\[
 a_n^{-1}U^-a_n\subseteq D,\qquad a_nU^+a_n^{-1}\subseteq D.
\]
The sign character is nontrivial on both $U^+$ and $U^-$. The first
inclusion shows that no $a_n$ intertwines it: on $U^-\subseteq
K\cap a_nKa_n^{-1}$, one of the two characters is nontrivial and the
other is trivial. Cartan decomposition gives the same conclusion for
every element outside $K$.

For completeness, the sign isotypic space under $K$ in
$\Pi=\cInd_K^G\operatorname{sgn}$ is therefore exactly its canonical
one-dimensional space supported on $K$, by the compact Mackey
decomposition. Any nonzero subrepresentation contains a vector whose
translate has nonzero value on that coset. Projecting to the sign
isotypic space gives a nonzero vector in the canonical copy, which
generates $\Pi$. Hence $\Pi$ is irreducible. The invariant unitary form
on compact induction has compactly supported matrix coefficients for
finite-support vectors; admissibility and irreducibility then imply
supercuspidality.

The second inclusion kills every $D$-fixed Mackey term away from $K$:
$U^+\subseteq K\cap a_n^{-1}Da_n$, whereas sign has no $U^+$-invariants.
Normality of $D$ in $K$ handles the compact factors. Consequently
$\Pi^D=\operatorname{sgn}$. Averaging a translate supported outside $K$
cannot create support on $K$; since every $D$-fixed vector is supported
on $K$, this average is zero. Thus the compressed action outside $K$
vanishes, and $\Pi^D$ is a $Q_1$-module. A finite-dimensional semisimple
complex algebra with
exactly one simple module, of dimension one, is $\C$, proving the claim.
\end{proof}

This example computes the exponent-one sector. For a general group or
level, the stopping bound $B_e$ can be larger than one. The full finite
algebra $Q^{\mathrm{sc}}_{e,\omega}$ then still parametrizes the
supercuspidals, but the preceding compact-induction argument does not
supply an inducing subgroup for every simple module.

\begin{thebibliography}{99}

\bibitem{Casselman}
W.~Casselman, \emph{Introduction to the theory of admissible representations
of $p$-adic reductive groups}, notes, 1 May 1995.
\url{https://www.math.utah.edu/~ptrapa/math-library/casselman/casselman-p-adic-book.pdf}.

\bibitem{DHKM}
J.-F.~Dat, D.~Helm, R.~Kurinczuk, and G.~Moss,
\emph{Finiteness for Hecke algebras of $p$-adic groups},
arXiv:2203.04929v2 (2022).
\url{https://arxiv.org/abs/2203.04929v2}.
Theorem numbers refer to this version.

\bibitem{HV}
G.~Henniart and M.-F.~Vign\'eras,
\emph{Representations of a reductive $p$-adic group in characteristic
distinct from $p$}, author's version, 30 October 2021.
\url{https://perso.imj-prg.fr/mariefrance-vigneras/wp-content/uploads/sites/74/2022/11/21-10-30.pdf}.

\bibitem{SS}
P.~Schneider and U.~Stuhler, \emph{Representation theory and sheaves on the
Bruhat--Tits building}, Publ. Math. IH\'ES \textbf{85} (1997), 97--191.
\url{https://www.numdam.org/item/PMIHES_1997__85__97_0/}.

\bibitem{Solleveld}
M.~Solleveld, \emph{Pseudo-reductive and quasi-reductive groups over
non-archimedean local fields}, author's version, 13 June 2018.
\url{https://www.math.ru.nl/~solleveld/quasiReductive.pdf}.
Theorem numbers refer to this version.

\end{thebibliography}

\end{document}

